\documentclass[11pt,a4paper]{article}

\usepackage[utf8]{inputenc}

\usepackage[T1]{fontenc}

\usepackage[polish]{babel}

\usepackage{lmodern}

\usepackage{geometry}

\geometry{margin=27mm}

\usepackage[hidelinks]{hyperref}

\title{Autonomia użytkownika, zaufanie i obciążenie opiekuna przy eskalacji alarmów AAL}

\author{Mateusz Pomianek}

\date{Szkic naukowy, wersja 1.0.0 | 28 listopada 2026}

\begin{document}

\maketitle

\begin{abstract}

Proponowana miara główna to różnica prawdopodobieństw trafnej eskalacji w H1.

\end{abstract}

\section{Wprowadzenie}

Model zaufania do automatyzacji uzasadnia osobny pomiar polegania na alarmie i deklarowanego zaufania (Lee \& See, 2004). Autonomia opisuje możliwość działania zgodnie z własną decyzją (Patrick et al., 2007). W scenariuszu jej wskaźnikiem będzie ocena kontroli nad udostępnieniem informacji przez osobę w roli użytkownika.

\section{Cel i pytania badawcze}

Hipotezy robocze wymagają zatwierdzenia przed rozpoczęciem właściwej weryfikacji.

\paragraph{H1.} W próbach z alarmem wymagającym interwencji jawne uzasadnienie zwiększy prawdopodobieństwo trafnej eskalacji w porównaniu z ogólnym komunikatem o tym samym priorytecie.

\paragraph{H2.} W scenariuszach niewymagających natychmiastowej reakcji możliwość samodzielnego potwierdzenia przekazania alarmu zwiększy ocenę kontroli użytkownika względem przekazania automatycznego.

\section{Materiał i metody}

Pierwszorzędowy wskaźnik: różnica prawdopodobieństw trafnej eskalacji w H1.

Wynik główny to różnica prawdopodobieństw trafnej eskalacji w H1, szacowana z modelu logistycznego z efektami uczestnika i scenariusza. Klucz poprawności ustala niezależny panel przed badaniem. Deklarowane obciążenie opiekuna po bloku alarmów i odsetek zbędnych eskalacji są wynikami dodatkowymi. Opiekunowie oceniają losowane scenariusze z niezależnie różnicowanym priorytetem i kosztem przerwania czynności. Możliwość potwierdzenia przekazania ocenia się w osobnym bloku z perspektywy użytkownika.

Przed badaniem z udziałem ludzi należy zatwierdzić wymagania etyczne, zasady zgody i ochronę danych.

\section{Plan analizy}

Zamrożenie estymandu, reguł oceny jakości i obsługi braków nastąpi przed dostępem do danych wynikowych. Raport obejmie oszacowanie kontrastu głównego z przedziałem niepewności i pełne mianowniki. Jeśli przewidziano hipotezy drugorzędowe, przed analizą należy ustalić sposób kontroli wielokrotności.

\section{Wyniki}

Nie zebrano danych empirycznych dla tej koncepcji.

\section{Dyskusja i ograniczenia}

Ograniczenia obejmują możliwe błędy pomiaru, dobór próby, alternatywne mechanizmy i ograniczoną możliwość uogólniania. Korelacja nie dowodzi przyczynowości; poprawność techniczna nie jest dowodem znaczenia klinicznego.

\section{Wnioski}

Hipotezy i sposób pomiaru stanowią podstawę oceny wykonalności oraz przygotowania właściwego protokołu.

\section*{Bibliografia}

{\setlength{\parindent}{0pt}
\setlength{\parskip}{5pt}
\hangindent=1.27cm
Lee, J. D., \& See, K. A. (2004). Trust in automation: Designing for appropriate reliance. \textit{Human Factors, 46}(1), 50--80. \url{https://doi.org/10.1518/hfes.46.1.50_30392}\par
\hangindent=1.27cm
Patrick, H., Knee, C. R., Canevello, A., \& Lonsbary, C. (2007). The role of need fulfillment in relationship functioning and well-being: A self-determination theory perspective. \textit{Journal of Personality and Social Psychology, 92}(3), 434--457. \url{https://doi.org/10.1037/0022-3514.92.3.434}\par
}

\end{document}
